\documentclass[9pt]{extarticle}
\usepackage[letterpaper,margin=0.45in]{geometry}
\usepackage{multicol}
\usepackage{amsmath,amssymb}
\usepackage{xcolor}
\usepackage{titlesec}
\usepackage{array,booktabs}
\usepackage{enumitem}
\usepackage{fancyhdr}

% --- Spacing ---
\setlength{\parindent}{0pt}
\setlength{\parskip}{2pt}
\setlength{\abovedisplayskip}{4pt}
\setlength{\belowdisplayskip}{4pt}
\setlength{\abovedisplayshortskip}{2pt}
\setlength{\belowdisplayshortskip}{2pt}
\setlength{\columnsep}{14pt}
\setlength{\columnseprule}{0.2pt}
\raggedcolumns

% --- Section formatting ---
\definecolor{gold}{HTML}{B8943E}
\definecolor{desc}{HTML}{444444}
\titleformat{\subsection}{\normalfont\bfseries\small\color{gold}}{}{0pt}{}
\titlespacing*{\subsection}{0pt}{5pt}{2pt}

% --- Header/Footer ---
\pagestyle{fancy}
\fancyhf{}
\renewcommand{\headrulewidth}{0pt}
\fancyfoot[L]{\footnotesize\textcopyright{} 2026 Ihsan S. All rights reserved.}
\fancyfoot[R]{\footnotesize\thepage/4}

% --- Compact list ---
\setlist{nosep,leftmargin=*,topsep=0pt}

% --- Shortcuts ---
\newcommand{\eqlabel}[1]{\textbf{\scriptsize #1}\;}
\newcommand{\desc}[1]{{\small\color{desc}#1}}

\begin{document}
\begin{multicols}{2}

% ============================================================
% QUANTUM MECHANICS (Lessons 1--2)
% ============================================================

\subsection{Physical Constants}
\renewcommand{\arraystretch}{1.25}
\begin{tabular}{@{}l@{\;\;}l@{}}
$h$ & $6.626\times10^{-34}$ J\,s \\
$\hbar = h/2\pi$ & $1.055\times10^{-34}$ J\,s \\
$hc$ & $1240$ eV\,nm \\
$m_e$ & $9.109\times10^{-31}$ kg \\
$e$ & $1.602\times10^{-19}$ C \\
$k_B$ & $1.381\times10^{-23}$ J/K = $8.617\times10^{-5}$ eV/K \\
$\varepsilon_0$ & $8.854\times10^{-12}$ F/m \\
$\mu_0$ & $4\pi\times10^{-7}$ N/A$^2$ \\
$c_0$ & $3.0\times10^{8}$ m/s \\
$N_A$ & $6.022\times10^{23}$ mol$^{-1}$ \\
$a_0$ & $0.0529$ nm (Bohr radius) \\
$e^2/(4\pi\varepsilon_0)$ & $1.44$ eV\,nm (Coulomb shortcut) \\
$k_BT|_{300\text{K}}$ & $0.02585$ eV \\
\end{tabular}
\renewcommand{\arraystretch}{1.0}

\subsection{Wave--Particle Duality}
\[E = hf = \hbar\omega\]
\desc{Energy of a photon; $f$ = frequency, $\omega = 2\pi f$}
\[p = \frac{h}{\lambda} = \frac{hf}{c}\]
\desc{Photon momentum}
\[\text{Photoelectric: } KE_{\max} = hf - \Phi\]
\desc{$\Phi$ = work function (min.\ energy to extract electron)}
\[\lambda = \frac{h}{p} = \frac{h}{\sqrt{2m_e\,eV}}\]
\desc{de Broglie wavelength; second form for electron through potential $V$}

\eqlabel{Wien} $\lambda_{\max}T = 2.898\times10^{-3}$ m\,K

\desc{Peak wavelength of blackbody radiation}

\eqlabel{Photon flux} $\dot{N} = \dfrac{P}{E_{ph}} = \dfrac{P\lambda}{hc}$

\desc{Photon rate from source of power $P$}

\eqlabel{Thermal $\lambda$} $\lambda_B = h/\sqrt{3mk_BT}$

\desc{Thermal de Broglie wavelength; quantum regime when $\lambda_B \gtrsim$ interparticle spacing}

\subsection{Schr\"odinger Equation}
\eqlabel{Momentum op.}
$\hat{p} = -i\hbar\dfrac{\partial}{\partial x}$

\eqlabel{Expectation}
$\langle A \rangle = \displaystyle\int \psi^* \hat{A}\,\psi\,dx$

\desc{Mean value of observable $A$; $\hat{A}$ is the quantum operator}

\eqlabel{Normalization}
$\displaystyle\int_{-\infty}^{\infty}|\psi|^2\,dx = 1$

\eqlabel{TDSE}
$i\hbar\dfrac{\partial}{\partial t}|\psi\rangle = \hat{H}|\psi\rangle$

\desc{Time-dependent Schr\"odinger equation}

\eqlabel{Hamiltonian}
$\hat{H} = -\dfrac{\hbar^2}{2m}\nabla^2 + V(\mathbf{r})$

\desc{Total energy operator = kinetic + potential}

\eqlabel{TISE}
$\dfrac{d^2\psi}{dx^2} + \dfrac{2m_e}{\hbar^2}(E - V)\psi = 0$

\desc{Time-independent form; solve for stationary states}

\subsection{Infinite Potential Well}
\desc{Particle confined to $0 < x < a$, $V = \infty$ outside}
\[\psi_n(x) = \sqrt{\frac{2}{a}}\sin\!\left(\frac{n\pi x}{a}\right)\]
\[E_n = \frac{n^2\pi^2\hbar^2}{2m_ea^2} = \frac{n^2 h^2}{8m_ea^2}\]
\desc{$a$ = well width, $n = 1,2,3,\ldots$ (quantum number)}
\[|\psi_n|^2 = \frac{2}{a}\sin^2\!\left(\frac{n\pi x}{a}\right)\]
\desc{Probability density; peaks at antinodes}
\[\Delta E_n = E_{n+1} - E_n = \frac{(2n+1)h^2}{8m_ea^2}\]

\eqlabel{3D box}
$E_{n_1n_2n_3} = \dfrac{h^2}{8m_e}\!\left(\dfrac{n_1^2}{a^2}+\dfrac{n_2^2}{b^2}+\dfrac{n_3^2}{c^2}\right)$

\desc{$a,b,c$ = box dimensions; degeneracy when $a=b=c$}

\subsection{Finite Potential Well}
$k^2 = \dfrac{2m_eE}{\hbar^2},\quad \alpha^2 = \dfrac{2m_e(V_0-E)}{\hbar^2}$

\desc{$k$ = wavevector inside well, $\alpha$ = decay constant outside}

\eqlabel{Even} $\alpha = k\tan\!\left(\dfrac{ka}{2}\right)$

\eqlabel{Odd} $\alpha = -k\cot\!\left(\dfrac{ka}{2}\right)$

\desc{Transcendental equations; solve graphically for bound state energies}

\eqlabel{Penetration} $\delta = \dfrac{1}{\alpha} = \dfrac{\hbar}{\sqrt{2m_e(V_0 - E)}}$

\desc{Depth of wavefunction tail into classically forbidden region}

\eqlabel{Bound states} $N \leq \left\lfloor \tfrac{1}{2} + \tfrac{a}{\pi}\sqrt{\tfrac{2m_eV_0}{\hbar^2}}\right\rfloor$

\subsection{Quantum Tunneling}
\[T = \frac{1}{1 + D\,\sinh^2(\alpha a)}\]
\[D = \frac{V_0^2}{4E(V_0-E)},\quad \alpha = \frac{\sqrt{2m_e(V_0-E)}}{\hbar}\]
\desc{$T$ = transmission probability through barrier of height $V_0$, width $a$}

\eqlabel{Thick barrier} $T \approx T_0\,e^{-2\alpha a},\; T_0 = \dfrac{16E(V_0-E)}{V_0^2}$

\desc{Valid when $\alpha a \gg 1$ (opaque barrier)}

\eqlabel{STM current} $I \propto V_b\,e^{-2\alpha d},\; \alpha = \dfrac{\sqrt{2m_e\phi}}{\hbar}$

\desc{Scanning tunneling microscope; $d$ = tip--sample gap, $\phi$ = avg.\ work function}

\subsection{Uncertainty Principle}
\[\Delta x\cdot\Delta p \geq \frac{\hbar}{2},\qquad \Delta E\cdot\Delta t \geq \frac{\hbar}{2}\]
\desc{Fundamental limits on simultaneous measurement precision}

\subsection{Hydrogen Atom}
\[V(r) = \frac{-Ze^2}{4\pi\varepsilon_0 r}\]
\desc{Coulomb potential; $Z$ = atomic number (1 for H)}
\[E_n = \frac{-13.6\text{ eV}}{n^2}\]
\desc{$n$ = principal quantum number; ground state $n=1$}
\[L = \hbar\sqrt{\ell(\ell+1)},\quad L_z = m_\ell\hbar\]
\desc{Orbital angular momentum; $\ell = 0\ldots n-1$}
\[S = \hbar\sqrt{s(s+1)} = \frac{\sqrt{3}}{2}\hbar,\quad S_z = m_s\hbar\]
\desc{Spin; $s = 1/2$ for electrons, $m_s = \pm 1/2$}

Quantum numbers: $n=1,2,\ldots;\; \ell=0\ldots n\!-\!1;\; m_\ell=-\ell\ldots+\ell;\; m_s=\pm\tfrac{1}{2}$

Degeneracy: $n^2$ orbitals, $2n^2$ states per shell

\eqlabel{Selection rules} $\Delta\ell=\pm1,\;\Delta m_\ell=0,\pm1$

\desc{Allowed electric dipole transitions}

\eqlabel{Virial} $\langle T \rangle = -E,\quad \langle V \rangle = 2E$

\desc{Avg.\ kinetic and potential energy from total $E$}

\eqlabel{Ionization} $I = -E_n = 13.6/n^2$ eV

\eqlabel{Bohr radius} $a_0 = \dfrac{4\pi\varepsilon_0\hbar^2}{m_e e^2} = 0.0529$ nm

% ============================================================
% SOLIDS (Lessons 3--4)
% ============================================================

\subsection{Interatomic Bonding}
\[E(r) = -\frac{A}{r} + \frac{B}{r^m}\]
\desc{Pair potential: attractive ($-A/r$) + repulsive ($B/r^m$, $m \approx 6$--$12$)}

\eqlabel{Ionic} $A = \dfrac{Me^2}{4\pi\varepsilon_0}$

\desc{$M$ = Madelung constant (depends on crystal geometry)}

\[E_{\text{lattice}}(r) = -\frac{Me^2}{4\pi\varepsilon_0\,r} + \frac{B}{r^m}\]
\[E_{\text{bond}} = \frac{A}{r_0}\left(1 - \frac{1}{m}\right)\]
\desc{Cohesive energy at equilibrium separation $r_0$}

\eqlabel{Equilibrium} $\left.\dfrac{dE}{dr}\right|_{r_0}=0$

$r_0 = \left(\dfrac{mB\cdot 4\pi\varepsilon_0}{e^2 M}\right)^{1/(m-1)}$

\subsection{Crystal Structures}
\renewcommand{\arraystretch}{1.15}
\begin{tabular}{@{}l c c c l@{}}
\toprule
& \small atoms & \small CN & \small APF & \small $r(a)$ \\
\midrule
SC & 1 & 6 & 0.52 & $a/2$ \\
BCC & 2 & 8 & 0.68 & $a\sqrt{3}/4$ \\
FCC & 4 & 12 & 0.74 & $a\sqrt{2}/4$ \\
DC & 8 & 4 & 0.34 & $a\sqrt{3}/8$ \\
\bottomrule
\end{tabular}
\renewcommand{\arraystretch}{1.0}

\desc{CN = coordination number, APF = atomic packing fraction}

\smallskip
\eqlabel{Density} $\rho = \dfrac{n_c\,M_a}{N_A\,a^3}$

\desc{$n_c$ = atoms/cell, $M_a$ = molar mass, $a$ = lattice constant}

\eqlabel{NaCl} $M = 1.748$\quad \eqlabel{CsCl} $M = 1.763$

\eqlabel{Surface} $\sigma_{100} = \dfrac{2}{a^2}$, \; $\sigma_{110} = \dfrac{2\sqrt{2}}{a^2}$, \; $\sigma_{111} = \dfrac{4}{\sqrt{3}\,a^2}$

\desc{Planar atomic density for diamond cubic (Si) [atoms/area]}

\subsection{X-ray Diffraction}
\[d_{hkl} = \frac{a}{\sqrt{h^2+k^2+l^2}}\]
\desc{Interplanar spacing for cubic crystals}
\[2d\sin\theta = n\lambda \quad\text{(Bragg's law)}\]
\desc{$\theta$ = diffraction angle, $n$ = order, $\lambda$ = X-ray wavelength}

Systematic absences:

BCC: $h\!+\!k\!+\!l$ = even \quad FCC: $h,k,l$ all odd or all even

\subsection{Crystal Defects}
\[n_v = N\exp\!\left(-\frac{E_v}{k_BT}\right)\]
\desc{$n_v$ = vacancy concentration, $E_v$ = vacancy formation energy}

\subsection{Stoichiometry}
$w_A = \dfrac{x\,M_A}{x\,M_A + y\,M_B}$, \quad $x_i = \dfrac{n_i}{\sum_j n_j}$

\desc{$w_A$ = weight fraction, $x_i$ = mole fraction}

\subsection{Band Theory}
\[E_g = E_c - E_v\]
\desc{$E_g$ = band gap; $E_c$ = conduction band edge; $E_v$ = valence band edge}

\eqlabel{Photoemission}
Metal: $hf > \Phi$, \; $KE = hf - \Phi$

Semiconductor: $hf > E_g + \chi$

\desc{$\chi$ = electron affinity ($E_{\text{vac}} - E_c$); $\Phi$ = work function ($E_{\text{vac}} - E_F$)}

\subsection{E--k Dispersion \& Effective Mass}
\[E = \frac{\hbar^2 k^2}{2m_e^*}\]
\desc{Parabolic approximation near band extrema}
\[\frac{1}{m_e^*} = \frac{1}{\hbar^2}\frac{d^2E}{dk^2}\]
\desc{Effective mass from curvature of $E$-$k$ band}

\subsection{Density of States}
\[g(E) = 8\pi\sqrt{2}\left(\frac{m_e}{h^2}\right)^{3/2}E^{1/2}\]
\desc{States per unit energy per unit volume (3D free electron)}
\[S_v(E') = \frac{\pi(8m_eE')^{3/2}}{3h^3}\]
\desc{Cumulative states up to energy $E'$}

\subsection{Fermi--Dirac Statistics}
\[f(E) = \frac{1}{1+\exp\!\left(\dfrac{E-E_F}{k_BT}\right)}\]
\desc{Probability that state at energy $E$ is occupied; $E_F$ = Fermi energy}

\eqlabel{Boltzmann approx.} $f \approx e^{-(E-E_F)/(k_BT)}$ for $E-E_F \gg k_BT$

\eqlabel{Boltzmann ratio} $\dfrac{N_2}{N_1} = \exp\!\left(-\dfrac{E_2-E_1}{k_BT}\right)$

\[E_{F0} = \frac{h^2}{8m_e}\left(\frac{3n}{\pi}\right)^{2/3}\]
\desc{Fermi energy at $T=0$; $n$ = free electron concentration}
\[E_F(T) = E_{F0}\!\left[1 - \frac{\pi^2}{12}\!\left(\frac{k_BT}{E_{F0}}\right)^{\!2}\right]\]
\[\bar{E} = \frac{3}{5}E_{F0}\]
\desc{Average electron energy at $T = 0$}
\[\langle E \rangle(T) = \tfrac{3}{5}E_{F0}\!\left[1 + \tfrac{5\pi^2}{12}\!\left(\tfrac{k_BT}{E_{F0}}\right)^{\!2}\right]\]
\desc{Temperature-corrected average energy}

\eqlabel{Fermi vel.} $v_F = \sqrt{2E_F/m_e}$

\eqlabel{Thermal vel.} $v_{\text{rms}} = \sqrt{3k_BT/m_e}$

\eqlabel{Inverse} $n = \dfrac{\pi}{3}\!\left(\dfrac{8m_eE_{F0}}{h^2}\right)^{3/2}$

\desc{Electron concentration from Fermi energy (inverted $E_{F0}$ formula)}

\eqlabel{2D DOS} $g_{2D} = \dfrac{4\pi m_e}{h^2}$ \desc{(energy-independent, per unit area)}

% ============================================================
% CONDUCTION \& OPTICS (Lesson 5)
% ============================================================

\subsection{Drude Model \& Conduction}
\[v_F = \sqrt{\frac{2E_F}{m_e}},\qquad l = v_F\tau\]
\desc{$v_F$ = Fermi velocity, $l$ = mean free path, $\tau$ = mean free time}
\[v_{dx} = \frac{eE_x}{m_e}\tau = \mu E_x\]
\desc{$v_{dx}$ = drift velocity in applied field $E_x$}
\[\mu = \frac{e\tau}{m_e}\]
\desc{$\mu$ = electron mobility [m$^2$/(V\,s)]}
\[J_x = en\,v_{dx} = \sigma E_x\]
\desc{$J_x$ = current density [A/m$^2$]}
\[\sigma = \frac{ne^2\tau}{m_e} = ne\mu\]
\desc{$\sigma$ = electrical conductivity [S/m]}
\[R = \frac{\rho\,L}{A},\qquad \rho = \frac{1}{\sigma}\]
\desc{$R$ = resistance, $\rho$ = resistivity, $L$ = length, $A$ = cross-section}

\eqlabel{Band model} $\sigma = \tfrac{1}{3}e^2v_F^2\tau\,g(E_F)$

\desc{Quantum conductivity from Fermi surface scattering}

\eqlabel{Electron conc.} $n = \dfrac{Z\cdot d\cdot N_A}{M_{at}}$

\desc{$Z$ = valence electrons, $d$ = mass density, $M_{at}$ = atomic mass}

\eqlabel{Mean free time} $\tau = \dfrac{m_e}{ne^2\rho}$

\desc{Extract $\tau$ from measured resistivity}

\subsection{Resistivity Temperature Dependence}
\[\rho(T) \approx \rho_0[1 + \alpha_0(T - T_0)]\]
\desc{Linear approximation; $\alpha_0$ = temperature coefficient of resistivity}

\eqlabel{Matthiessen} $\rho = \rho_T + \rho_I$

\desc{Total = thermal scattering + impurity scattering (additive)}

\subsection{Thermoelectrics}
\[S = \frac{dV}{dT} \quad\text{(Seebeck coeff.)}\]
\desc{Voltage generated per unit temperature difference [V/K]}
\[V_{AB} = \int_{T_0}^{T}(S_A - S_B)\,dT\]
\desc{Thermocouple voltage between materials $A$ and $B$}

\eqlabel{Contact potential} $\Delta V = \dfrac{\Phi_1 - \Phi_2}{e}$

\desc{Built-in voltage from work function difference at junction}

\subsection{Maxwell's Equations}
\begin{align*}
\nabla\cdot\vec{D} &= \rho_{\text{free}} & \nabla\cdot\vec{B} &= 0 \\
\nabla\times\vec{E} &= -\frac{\partial\vec{B}}{\partial t} & \nabla\times\vec{H} &= \vec{J} + \frac{\partial\vec{D}}{\partial t}
\end{align*}

\subsection{EM Waves \& Optics}
\[\nabla^2\mathbf{E} - \frac{1}{c^2}\frac{\partial^2\mathbf{E}}{\partial t^2} = 0\]
\desc{Wave equation in a linear medium}
\[c_0 = \frac{1}{\sqrt{\varepsilon_0\mu_0}},\quad n = \sqrt{\varepsilon_r},\quad c = \frac{c_0}{n}\]
\desc{$n$ = refractive index, $\varepsilon_r$ = relative permittivity}

\eqlabel{Intensity} $I = \tfrac{1}{2}c\varepsilon_0 E_0^2$

\desc{Time-averaged power per unit area for plane wave}

\eqlabel{Sellmeier} $n^2(\lambda) = 1 + \displaystyle\sum_i\frac{B_i\lambda^2}{\lambda^2 - C_i}$

\desc{Empirical dispersion; $B_i$, $C_i$ = material coefficients}

\[v_g = \frac{c}{N_g},\quad N_g = n - \lambda\frac{dn}{d\lambda}\]
\desc{$v_g$ = group velocity, $N_g$ = group index}

\eqlabel{Plane wave} $E(z,t) = E_0\cos(kz - \omega t + \phi_0)$

$k = \dfrac{2\pi}{\lambda},\quad \omega = 2\pi f,\quad v = f\lambda = \dfrac{\omega}{k}$

\subsection{Snell's Law \& Fresnel}
\[n_1\sin\theta_i = n_2\sin\theta_t\]
\desc{$\theta_i$ = angle of incidence, $\theta_t$ = angle of transmission}
\[R_s = \left(\frac{n_1\cos\theta_i - n_2\cos\theta_t}{n_1\cos\theta_i + n_2\cos\theta_t}\right)^{\!2}\]
\desc{$s$-polarized reflectance ($\perp$ to plane of incidence)}
\[R_p = \left(\frac{n_2\cos\theta_i - n_1\cos\theta_t}{n_2\cos\theta_i + n_1\cos\theta_t}\right)^{\!2}\]
\desc{$p$-polarized reflectance ($\parallel$ to plane of incidence)}

\eqlabel{Normal inc.} $R = \left(\dfrac{n_1-n_2}{n_1+n_2}\right)^{\!2}$

\[\tan\theta_B = \frac{n_2}{n_1} \quad\text{(Brewster)}\]
\desc{Angle where $R_p = 0$; reflected light is purely $s$-polarized}
\[\theta_c = \arcsin\!\left(\frac{n_2}{n_1}\right) \quad\text{(Critical)}\]
\desc{Total internal reflection for $\theta_i > \theta_c$ (requires $n_1 > n_2$)}

\eqlabel{Fiber NA} $\text{NA} = \sin\theta_{\text{in,max}} = \sqrt{n_1^2 - n_2^2}$

\desc{Numerical aperture; max acceptance angle for waveguide}

\subsection{Absorption \& Attenuation}
\[I(z) = I_0\,e^{-\alpha z}\]
\desc{Beer--Lambert law; $\alpha$ = absorption coefficient [Np/m]}
\[\alpha = \frac{\pi\varepsilon_r''}{\lambda\sqrt{\varepsilon_r'}}\]
\[\frac{P_{\text{out}}}{P_{\text{in}}} = e^{-2\alpha z}\]
\desc{Power attenuation uses $2\alpha$ (intensity $\propto E^2$)}

\subsection{Complex Refractive Index}
$N = n + j\kappa$

\desc{$n$ = refractive index, $\kappa$ = extinction coefficient}

\eqlabel{Absorption} $\alpha = \dfrac{2\omega\kappa}{c} = \dfrac{4\pi\kappa}{\lambda}$

% ============================================================
% DIELECTRICS (Lesson 6)
% ============================================================

\subsection{Electric Polarization}
\[\mathbf{D} = \varepsilon_0\mathbf{E} + \mathbf{P} = \varepsilon_r\varepsilon_0\mathbf{E}\]
\desc{$\mathbf{D}$ = displacement field, $\mathbf{P}$ = polarization}
\[\mathbf{P} = \chi_e\varepsilon_0\mathbf{E},\qquad \varepsilon_r = 1 + \chi_e\]
\desc{$\chi_e$ = electric susceptibility}
$\mathbf{P} = N\cdot\mathbf{p}_{\text{avg}}$, \quad $\mathbf{p} = \alpha_e\,\mathbf{E}_{\text{loc}}$

\desc{$N$ = number density, $\alpha_e$ = polarizability}
\[E_{\text{loc}} = E + \frac{P}{3\varepsilon_0}\]
\desc{Lorentz local field at atomic site $\neq$ applied field}

\subsection{Clausius--Mossotti}
\[\frac{\varepsilon_r - 1}{\varepsilon_r + 2} = \frac{N\alpha_e}{3\varepsilon_0}\]
\desc{Links macroscopic $\varepsilon_r$ to microscopic polarizability $\alpha_e$}

Rearranged: $\alpha_{\text{total}} = \dfrac{3\varepsilon_0(\varepsilon_r - 1)}{N(\varepsilon_r + 2)}$

\subsection{Polarizability Mechanisms}
\[\alpha_{\text{total}} = \alpha_e + \alpha_i + \alpha_d\]
\desc{Electronic + ionic + orientational (dipolar)}
\[\alpha_e = 4\pi\varepsilon_0 R^3 \quad\text{(electronic)}\]
\desc{$R$ = atomic/ionic radius; dominates at all frequencies}
\[\alpha_d = \frac{p^2}{3k_BT} \quad\text{(orientational)}\]
\desc{$p$ = permanent dipole moment; active only at low $\omega$}

\eqlabel{Langevin} $L(x) = \coth(x) - \dfrac{1}{x},\; x = \dfrac{pE}{kT}$

\desc{Exact orientational polarization; $\alpha_d$ is the weak-field limit}

\subsection{Optical Permittivity}
$\alpha_{\text{opt}} = \alpha_e(\text{cation}) + \alpha_e(\text{anion})$

$\dfrac{\varepsilon_{rop}-1}{\varepsilon_{rop}+2} = \dfrac{N\alpha_{\text{opt}}}{3\varepsilon_0}$

\desc{At optical frequencies only electronic polarizability contributes}

\subsection{Debye Relaxation}
\[\varepsilon_r(\omega) = \varepsilon_{r\infty} + \frac{\varepsilon_{rs}-\varepsilon_{r\infty}}{1+j\omega\tau}\]
\desc{$\varepsilon_{rs}$ = static (DC), $\varepsilon_{r\infty}$ = high-freq., $\tau$ = relaxation time}
\[\varepsilon_r' = \varepsilon_{r\infty} + \frac{\varepsilon_{rs}-\varepsilon_{r\infty}}{1+\omega^2\tau^2}, \quad \varepsilon_r'' = \frac{(\varepsilon_{rs}-\varepsilon_{r\infty})\omega\tau}{1+\omega^2\tau^2}\]
\desc{$\varepsilon_r'$ = storage (real), $\varepsilon_r''$ = loss (imaginary)}
\[\tan\delta = \frac{\varepsilon_r''}{\varepsilon_r'}\]
\desc{Loss tangent; ratio of energy dissipated to stored per cycle}

\desc{Peak of $\varepsilon_r''$ at $\omega\tau = 1$: \; $\varepsilon_{r,\text{peak}}'' = (\varepsilon_{rs}-\varepsilon_{r\infty})/2$}

\eqlabel{Complex $\varepsilon$} $\varepsilon_r(\omega) = \varepsilon_r'(\omega) - j\varepsilon_r''(\omega)$

\eqlabel{Complex $k$} $\tilde{k} = \dfrac{\omega}{c}\sqrt{\varepsilon_r} = k' + ik''$

\subsection{Wave in Lossy Medium}
\[E(z) = E_0\,e^{-k''z}\,e^{i(k'z-\omega t)}\]
\desc{$k''$ attenuates amplitude; $k'$ determines phase velocity}

\subsection{Gauss's Law for D}
$\nabla\cdot\mathbf{D} = \rho_{\text{free}}$, \quad $\displaystyle\oint\mathbf{D}\cdot d\mathbf{A} = Q_{\text{free,enclosed}}$

\eqlabel{Boundary} $D_{n1} - D_{n2} = \sigma_f$

\desc{$\sigma_f$ = free surface charge density at interface}

\subsection{Capacitors \& Energy}
\[C = \frac{\varepsilon_r\varepsilon_0 A}{d}\]
\desc{$A$ = plate area, $d$ = plate separation}
\[u = \frac{1}{2}\varepsilon_r\varepsilon_0 E^2 = \frac{1}{2}\mathbf{D}\cdot\mathbf{E}\]
\desc{Energy density [J/m$^3$]; total $U = \tfrac{1}{2}CV^2$}

\eqlabel{Parallel} $C_{\text{tot}} = C_1 + C_2 + \cdots$
\quad
\eqlabel{Series} $1/C_{\text{tot}} = 1/C_1 + 1/C_2 + \cdots$

\eqlabel{Multilayer} $C = \dfrac{\varepsilon_0 A}{d_1/\varepsilon_{r1} + d_2/\varepsilon_{r2}}$

\desc{Two dielectric layers in series}

\subsection{Piezoelectricity}
\[P = d\cdot T \;\text{(direct)},\quad S = d\cdot E \;\text{(converse)}\]
\desc{$d$ = piezoelectric coefficient, $T$ = stress, $S$ = strain}
\[f = \frac{v}{2L} \quad\text{(resonance)},\quad Q = \frac{f_{\text{res}}}{\Delta f}\]
\desc{$v$ = acoustic velocity in crystal, $L$ = thickness, $Q$ = quality factor}

\subsection{Standing Waves \& Interference}
\[E = 2E_0\cos(kz)\cos(\omega t)\]
\desc{Superposition of counter-propagating waves; antinode spacing $\lambda/2$}

\eqlabel{Double slit} $\Delta = d\sin\theta$

\desc{Path difference; constructive $\Delta = n\lambda$, destructive $\Delta = (n+\tfrac{1}{2})\lambda$}

\subsection{Superposition (Infinite Well)}
\[\Psi(x,t) = c_1\varphi_1 e^{-iE_1t/\hbar} + c_2\varphi_2 e^{-iE_2t/\hbar}\]
\[|\Psi|^2 = c_1^2\varphi_1^2 + c_2^2\varphi_2^2 + 2c_1c_2\varphi_1\varphi_2\cos\!\left(\frac{(E_2-E_1)t}{\hbar}\right)\]
\desc{Oscillation frequency: $\omega_{21} = (E_2 - E_1)/\hbar$}

\subsection{Material Properties}
\renewcommand{\arraystretch}{1.15}
\begin{tabular}{@{}l c c@{}}
\toprule
Material & $a$ (nm) & $E_g$ (eV) \\
\midrule
Si & 0.543 & 1.12 \\
Ge & 0.566 & 0.66 \\
GaAs & 0.565 & 1.42 \\
Diamond & 0.357 & 5.5 \\
\bottomrule
\end{tabular}
\renewcommand{\arraystretch}{1.0}

\smallskip
Si: $n_{\text{Si}} = 4.99\times10^{22}$ cm$^{-3}$, $\rho = 2.33$ g/cm$^3$

$n_{\text{Si}}$(refractive) $= 3.48$, \; $n_{\text{SiO}_2} = 1.44$

\subsection{Spectroscopic Notation}
$\ell = 0,1,2,3 \to s,p,d,f$

Shell capacity: $2(2\ell+1)$ electrons per subshell

\subsection{Useful Identities \& Conversions}
$1\text{ eV} = 1.602\times10^{-19}$ J, \; $1\text{ nm} = 10^{-9}$ m, \; $1\text{ \AA} = 0.1\text{ nm}$

$E = hc/\lambda$: at $\lambda = 1240\text{ nm}$, $E = 1\text{ eV}$

\eqlabel{Euler} $e^{i\theta} = \cos\theta + i\sin\theta$

$\sinh(x) = (e^x - e^{-x})/2$, \; $\cosh(x) = (e^x + e^{-x})/2$

\eqlabel{Stopping} $eV_{\text{stop}} = KE_{\max} = hf - \Phi$

$\hbar c = 197.3$ eV\,nm, \; $m_ec^2 = 0.511$ MeV

\eqlabel{Dispersion} $E = \hbar\omega$, $p = \hbar k$

\eqlabel{Free electron} $\psi = Ae^{jkx}$, $k = \sqrt{2m_eE}/\hbar$

$R + T = 1$ (reflection + transmission = 1)

$N = n_c/a^3$ (number density from unit cell)

\end{multicols}
\end{document}
